\section{The finite problem and its geometric meaning}
\label{sec:model}

\subsection{Rows, columns, and the rooted dual graph}

Let $D$ be a finite connected multigraph with distinguished root $\ast$.
Its non-root vertices represent three-handles.  Each row $r\in R$ of a
nonnegative integer matrix $A=(a_{rc})$ also labels an edge of $D$.
Rows not eligible for a top cancellation are represented by root loops;
the software permits an equivalent \code{None} edge record.  Parallel
edges and loops remain distinct occurrences.

Columns represent retained one-handle cocores after a fixed lower
zero/one-handle forest has been contracted.  An entry $a_{rc}$ is an
\emph{actual unsigned number of crossings} of an attaching two-handle
curve with that cocore.  It is not the signed exponent sum of a word.
For the raw cancellation decision, all values above one may be saturated
at two.  This preserves whether a restriction has a column sum exactly
one.  Saturation must not be used for an unrelated parity computation.

Deleting the rows of a spanning tree $T$ models the source-certified
top cancellations.  A private-column peel subsequently removes a row
$r$ when some column has exactly one total remaining occurrence, located
on that row.  It removes that pivot column as well; all other entries of
the remaining rows stay unchanged.  A column that has already served as
a private pivot cannot later matter to an active row, because it had
no other live owner when removed.

\begin{definition}[Forest-compatible raw completion]
A pair $(D,A)$ admits a forest-compatible raw completion when some spanning
tree $T$ of $D$ leaves a row restriction $A|_{R\setminus T}$ that can be
fully removed by private-column peeling.
\end{definition}

This is the exact finite problem solved below.  Under the source-authentication
hypotheses of the preceding geometry compiler, success gives a valid
topological cancellation sequence.  Without those hypotheses, it is only
a graph/matrix theorem.

\subsection{Why stopping sets are the right obstruction}

\begin{definition}[Stopping set]
A nonempty row set $S$ is stopping if, for each column $c$,
$\sum_{r\in S}a_{rc}$ is either zero or at least two.
\end{definition}

If all rows of a stopping set remain, no first private-column peel can
remove a row from that set.  Conversely, a nonempty residual of exhaustive
column-leaf peeling is stopping.  The predecessor establishes the stronger
fact that allowing arbitrary positive singleton substitutions does not
enlarge the class of matrices admitting complete raw cancellation
\cite{IncomingPositive2026,IncomingForests2026}.  We use that exact
completeness result to focus on no-fill peeling.  We do not introduce
negative coefficients, free reduction, or unverified geometric isotopy.

The obstruction is restricted to the chosen calculus.  A stopping set
does not prove that a manifold is not a handlebody, and certainly does
not prove that a knot is nontrivial.  Conversely, some native stopping
examples really are non-handlebodies, so not every refusal should be
treated as a positive instance awaiting a clever cancellation order.

\subsection{The existing polynomial regime}

Call a column \emph{inert} if it has no entry equal to one.  Such a column
cannot ever be used by a private-column peel, even after row deletion.
Ignore it for the finite decision while retaining any corresponding
geometric handle count.

For a useful column supported on a single row $u$, put $u$ in a seed set
$Z$.  For a useful column supported on two rows $u,v$, add $u\to v$ when
the entry on $v$ is one, and add $v\to u$ when the entry on $u$ is one.
The other entry may exceed one.  These directed arcs form the ordinary
propagation graph $P$.  From any deleted row set $Q$, ordinary peeling
removes exactly the rows reachable from $Z\cup Q$.

If all useful columns are ordinary, every source strongly connected
component of $P$ without a seed must be hit by the tree.  These components
are disjoint.  Independent representatives from them can be selected by
intersecting the graphic matroid of $D$ with a partition matroid and then
extending the selected forest to a spanning tree.  This is the predecessor's
polynomial theorem, based on the classical independent-representatives
and matroid-intersection results
\cite{Rado1942,Edmonds1979,Lawler1975,IncomingForests2026}.

\subsection{Why native triple columns are structurally natural}

An original native one-handle cocore comes from a complementary region of
a triangular source face cut by normal arcs.  Such a region meets each of
the triangle's three sides in at most one gap.  Each unsigned two-handle
occurrence corresponds to an actual boundary gap.  Therefore the total
number of occurrences on the native cocore is at most three.  Lower-forest
contraction deletes columns, and top cancellation deletes rows; neither
increases this bound.  The preceding normalization report proves this
with source-occurrence conventions \cite{IncomingForests2026}.

The possible positive-entry multisets are consequently
\[
(1),\quad(2),\quad(3),\quad(1,1),\quad(1,2),\quad(1,1,1).
\]
Only the last type is outside ordinary propagation.  It supplies three
rules, each with two premises: $\{u,v\}\to w$, $\{u,w\}\to v$, and
$\{v,w\}\to u$.  This explains the focus on triple-unit columns.  The
stronger Boolean theorem also applies to wider exceptional columns, with
the union of their incident rows as its parameter.

The native bound does not automatically survive arbitrary interface
augmentation, handle slides, or a change of representation.  A native
adapter must measure the actual exported matrix and check the stated
support conditions.  Our source audit does so, and the generic Boolean
solver treats wider columns without pretending that they satisfy the
three-occurrence theorem.

\subsection{Size and success conventions}

The input size $N$ includes explicit row and edge records and all integer
bit lengths.  The number of rows is $m$, the graphic rank is
$r=|V(D)|-1$, the number of exceptional columns is $k$, and their incident-row
union is $U$, with $t=|U|$.  The source graph is checked for connectivity
before arrays indexed by its vertex namespace are allocated.

A positive certificate consists of a spanning tree and a complete sequence
of row/column pivot pairs.  The unchanged predecessor verifier reads the
original sparse matrix, checks the tree, and checks each actual column
sum.  Randomized zero tests and resource limits are inconclusive.  The
deterministic fallback can conclude exhaustive finite-model infeasibility;
the present implementation does not export a compact negative certificate
for that entire search.  This distinction is recorded in its output schema.
